\subsection{Méthodologie}
\label{sec:methodologie_sans_retrait}

Les simulations reposent sur les hypothèses suivantes :
\begin{itemize}
    \item un investisseur de moins de 70 ans disposant d'un capital (variable) initial à investir ;
    \item deux situations patrimoniales sont étudiées :
    \begin{enumerate}
        \item un patrimoine hors CTO et assurance-vie de 1 million d'euros ;
        \item un patrimoine hors CTO et assurance-vie de 300 000 euros;
    \end{enumerate}
    \item un horizon d'analyse de plusieurs décennies, afin d'intégrer l'effet des rendements et des frais de gestion;
    \item des hypothèses fiscales correspondant à la législation en vigueur en 2025.
\end{itemize}
Les figures représentent la différence d'héritage net relatif, définie comme
\begin{equation}
    \frac{H_{AV} - H_{CTO}}{B + C_{CTO}}
\end{equation}
en fonction de la durée de détention, des frais de gestion et du rendement. Les valeurs positives indiquent un avantage de l'assurance-vie, tandis que les valeurs négatives traduisent un avantage du CTO.
